\documentclass[11pt]{article}
\usepackage[margin=2.2cm]{geometry}
\usepackage{ctex}
\usepackage{amsmath,amssymb}
\usepackage{booktabs}
\usepackage{longtable}
\usepackage{array}
\providecommand{\cmrsideswitch}{}
\usepackage{xcolor}
\usepackage[colorlinks=true,linkcolor=blue,urlcolor=blue]{hyperref}
\renewcommand{\arraystretch}{1.25}
\setlength{\parskip}{0.4em}
\title{因子降维 质疑与结论}
\author{常笑涵}
\date{Semester 2, 2026}
\begin{document}
\maketitle
\tableofcontents
\newpage
\begin{quote}
\textbf{一句话}：我们把 24 个因子"降维"成少数几块，\textbf{降维这件事本身成立}（输出几乎不变、参数降 5 倍），\textbf{但挂在它后面的"复杂加权 + 收缩"没有增量}。

\textbf{日期}：2026-10-11　\textbf{性质}：阶段性汇报（对外沟通用）

\textbf{标注}：🟢【课内】= 课程讲义讲过 \textperiodcentered{} 🟡【半课内】= 讲义提过、中间推理由我们补全 \textperiodcentered{} 🔴【发散】= 课程外、业界通用做法
\end{quote}
\subsection{\texorpdfstring{零、先回答一个问题：我们是不是走极端了？}{零、先回答一个问题：我们是不是走极端了？}}
\textbf{没有。} 看起来像"全盘否定"，是因为\textbf{我们把四件事捆在一起叫"降维"}，而它们各自的表现完全不同：
{\small
\begin{longtable}{p{5.02cm}p{1.63cm}p{9.54cm}}
\toprule
捆在"降维"里的四件事 & 判定 & 关键证据 \\
\midrule
\textbf{(1) 减列数}（24 个因子 \ensuremath{\rightarrow} G 个因子块） & \textbf{[已完成]} \textbf{成立} & \ensuremath{\Sigma} 的参数 \textbf{300 \ensuremath{\rightarrow} 55}（G=10）；条件数 \textbf{2020 \ensuremath{\rightarrow} 72.6}；输出秩相关 \textbf{0.9+} \\
(2) 块内等权带来的\textbf{隐式重加权} & \textbf{[注意]} 方向不稳 & 等价于"按块大小反比加权"，不是压缩的功劳 \\
(3) \textbf{GLS 加权}（按相关结构配权重） & ❌ 无增量 & 配对检验不显著；某些设定下\textbf{显著更差} \\
(4) \textbf{收缩}（压住极端权重） & ❌ 无增量 & 5 折交叉验证选出 \ensuremath{\lambda} = 0，即"不需要收缩" \\
\bottomrule
\end{longtable}
}
\begin{quote}
\textbf{所以不是"降维没用"，是"降维有用，但挂在它后面的三件附加品没用"。} 结论是\textbf{收窄}，不是否定。
\end{quote}
\subsection{\texorpdfstring{一、这件事在做什么（大白话）}{一、这件事在做什么（大白话）}}
Alpha 策略的链路是：\textbf{因子 \ensuremath{\rightarrow} 综合打分 \ensuremath{\rightarrow} 优化器 \ensuremath{\rightarrow} 交易}。
因子的麻烦是\textbf{数量多、彼此重复}。比如"波动率""特质波动率""日内振幅"是三个名字、说的基本是同一件事 ------ 全塞进打分，等于\textbf{同一件事被重复投了三票}。
\textbf{降维的目的}：把 24 个因子压成少数几个"块"，\textbf{在结果基本不变的前提下，大幅减少下游要估的参数}。
\begin{quote}
\textbf{注意}：降维的\textbf{目的不是让结果更好}，而是\textbf{让结果差不多、但算得起}。这一点很关键，因为判据完全不同：该比的是"和未压缩版本差多少"，不是"跑不跑得赢单因子"。
\end{quote}
\subsection{\texorpdfstring{二、目前做完了什么}{二、目前做完了什么}}
{\small
\begin{longtable}{p{1.92cm}p{6.90cm}p{7.38cm}}
\toprule
阶段 & 做了什么 & 结论 \\
\midrule
\textbf{Step 0} & 诊断：因子池里到底有没有冗余 & 24 个因子的\textbf{有效维度只有 3.87}（大量冗余）\textbf{[已完成]} \\
\textbf{Step 1} & 扩池：12 \ensuremath{\rightarrow} \textbf{24 个因子} & 有显著性的因子从 5/12 升到 \textbf{17/24} \\
\textbf{Step 2} & 聚类 + 诚实基准对照 & 见下 \\
\textbf{Step 3} & \textbf{压缩曲线}（本次重点） & 见下 \\
\textbf{Step 4} & \textbf{F 检验 / 等价性检验}（按你们习惯的口径） & 见下 \\
\textbf{质疑流程} & 6 类对抗性检验 & 见下（部分仍在跑） \\
\bottomrule
\end{longtable}
}
\textbf{顺带修正了一个数据口径 bug}：旧的价值因子（市净率倒数、盈利收益率）\textbf{分母用错了}（拿流通市值当分母，量纲差一个股本因子）。改正后 \textbf{市净率倒数的 t 值从 1.78 升到 8.14} ------ 价值因子的信号一直都在，是被量纲错误压没了。
\subsection{\texorpdfstring{三、数据支持的结论（正面）}{三、数据支持的结论（正面）}}
\subsubsection{\texorpdfstring{3.1 输出高度一致}{3.1 输出高度一致}}
压缩到 \textbf{G=10} 时，压缩后的打分与未压缩打分的\textbf{日横截面秩相关 0.905（最差的一天 0.828）}。
\textbf{秩相关的 95\% 置信区间只有 \ensuremath{\pm}0.004 宽}，精度足够下结论 ------ 这是本次最稳的一条。
\subsubsection{\texorpdfstring{3.2 下游计算负担大幅下降}{3.2 下游计算负担大幅下降}}
{\small
\begin{longtable}{p{3.76cm}p{4.70cm}p{2.82cm}p{1.88cm}p{3.05cm}}
\toprule
G（块数） & 与未压缩输出的秩相关 & \ensuremath{\Sigma} 的参数个数 & 条件数 & 参数/有效样本 \\
\midrule
\textbf{24（未压缩）} & 1.000 & \textbf{300} & \textbf{2020} & \textbf{5.18} \\
12 & 0.928 & 78 & 95.4 & 1.35 \\
\textbf{10} & \textbf{0.905} & \textbf{55} & \textbf{72.6} & \textbf{0.95} \\
\textbf{6} & 0.888 & \textbf{21} & 41.2 & \textbf{0.36} \\
2 & 1.000 & 3 & 4.33 & 0.05 \\
\bottomrule
\end{longtable}
}
\begin{quote}
\textbf{"参数/有效样本"} = 要估的参数个数 \ensuremath{\div} 有效独立样本数。\textbf{大于 1 就意味着参数比数据还多} ------ 未压缩版本是 5.18，压缩到 G=10 才降到 0.95（够用），G=6 是 0.36（宽松）。
\end{quote}
\subsubsection{\texorpdfstring{3.3 非劣性检验通过（\ensuremath{\delta} = 10\%）}{3.3 非劣性检验通过（δ = 10）}}
在验证窗上，压缩版与未压缩版的\textbf{日 IC 之差}，95\% 置信区间落在 \ensuremath{\pm}10\% 可忽略带内：
{\small
\begin{longtable}{p{2.35cm}p{5.23cm}p{4.44cm}p{4.18cm}}
\toprule
G & \ensuremath{\Delta}IC（相对参考） & 95\% 置信区间 & \textbf{不差过 10\%？} \\
\midrule
15 & \ensuremath{-}12.2\% & [\ensuremath{-}17.2\%, \ensuremath{-}7.2\%] & ❌ \\
12 & +1.5\% & [\ensuremath{-}6.0\%, +9.0\%] & \textbf{[已完成]} \\
10 & +24.6\% & [+12.9\%, +36.3\%] & \textbf{[已完成]} \\
8 & \ensuremath{-}4.1\% & [\ensuremath{-}8.7\%, +0.6\%] & \textbf{[已完成]} \\
6 & +37.1\% & [+27.1\%, +47.1\%] & \textbf{[已完成]} \\
4 / 3 / 2 & +23.6\% / +26.3\% / 0\% & 均落在带内 & \textbf{[已完成]} \\
\bottomrule
\end{longtable}
}
\textbf{除了 G=15，全部通过。}
\subsubsection{\texorpdfstring{3.4 分块是稳定结构，不是偶然}{3.4 分块是稳定结构，不是偶然}}
把训练窗切成 4 个不重叠的子期，各自重新聚类，G=10 的分块\textbf{两两一致性 0.964 \textasciitilde{} 1.000}。
\subsubsection{\texorpdfstring{3.5 池子里不止一个效应}{3.5 池子里不止一个效应}}
剔除最强的"反彩票/低波动"一族（7 个因子）后，\textbf{剩下 17 个因子仍能保住 82\% 的预测力}（IC 0.1345 vs 全池 0.1649），而且它们之间\textbf{仍有可压缩的冗余}（有效维度 9.58 / 17）。
\subsection{\texorpdfstring{四、数据不支持的结论（负面）}{四、数据不支持的结论（负面）}}
\subsubsection{\texorpdfstring{4.1 GLS 加权没有增量}{4.1 GLS 加权没有增量}}
按"谁有效 + 谁和谁不重复"配置权重（GLS），配对检验显示\textbf{与简单基准在统计上分不出来}；在样本外的某些设定下\textbf{反而显著更差}（因为它给出负权重，在训练集上过拟合）。
\subsubsection{\texorpdfstring{4.2 收缩不需要}{4.2 收缩不需要}}
用训练窗内部的 5 折交叉验证选收缩强度，\textbf{选出的是 \ensuremath{\lambda} = 0，即"完全不需要收缩"}。这与上一轮结论一致：块层面的协方差本来就良态（10 个块、1217 天），不需要额外补救。
\subsubsection{\texorpdfstring{4.3 "按相关性聚类"这个判据本身不指向效果（\textbf{短窗冒烟，全量待跑}）}{4.3 "按相关性聚类"这个判据本身不指向效果（短窗冒烟，全量待跑）}}
{\small
\begin{longtable}{p{1.27cm}p{3.51cm}p{2.81cm}p{8.61cm}}
\toprule
G & 按相关性聚类的一致性 & 随机分组的一致性 & 关键相关 \\
\midrule
10 & \textbf{0.412} & \ensuremath{\approx} 0 & \textbf{随机分组中"一致性越高 \ensuremath{\rightarrow} 效果越差"（\ensuremath{-}0.398）} \\
8 & \textbf{0.368} & \ensuremath{\approx} 0 & +0.057 \\
\bottomrule
\end{longtable}
}
按相关性聚类确实把"相似"的因子归到一块（一致性 0.41 vs 随机 0），\textbf{但没有任何证据表明这转化为更好的结果}，某些设定下甚至是反向的。
\subsubsection{\texorpdfstring{4.4 损失不是单调的}{4.4 损失不是单调的}}
一个真实的"信息损失"应该随压缩加剧而\textbf{单调变差}。实测的 Spearman 相关是 \textbf{\ensuremath{-}0.571}（反号），说明波动主要由重加权驱动，不是信息在流失。
\subsection{\texorpdfstring{五、$\star$ 本次最重要的一个发现：压缩本身零损失}{五、★ 本次最重要的一个发现：压缩本身零损失}}
我们把块权重设成 \textbf{∝ 块大小} 之后，每个因子的权重是
\[
e_f=\frac{|block_g|}{J}\cdot\frac{\mathrm{sgn}_f}{|block_g|}=\frac{\mathrm{sgn}_f}{J}
\]
\textbf{这恒等于纯等权。} 因此：
\begin{quote}
\textbf{纯压缩不改动任何输出（逐日精确相等）。前面所有 \ensuremath{\Delta}IC 数字（\ensuremath{\pm}0.001 \textasciitilde{} \ensuremath{\pm}0.033）100\% 来自"块大小不均造成的隐式重加权"，与压缩本身无关。}
\end{quote}
\textbf{推论}：\textbf{任何"块大小相等"的分块（例如 24/6=4 个一组）给出的分数与等权完全相同。压缩只有通过"块大小不均"才对输出产生作用。}
这条同时说明：\textbf{要检验"压缩损失多少"，必须先把重加权拆出去，否则测的是混合物。}
\subsection{\texorpdfstring{六、一个必须澄清的比较口径}{六、一个必须澄清的比较口径}}
"压缩损失"要跟谁比？这里有两个完全不同的基准：
{\small
\begin{longtable}{p{2.62cm}p{5.40cm}p{8.18cm}}
\toprule
基准 & 是什么 & 压缩要追平它的难度 \\
\midrule
\textbf{等权 24 因子} & 每个因子各 1/24 & 容易（3.3 节基本都通过） \\
\textbf{最优未压缩} & 24 个因子\textbf{各配一个估出来的权重} & \textbf{更难}（对着它，\ensuremath{\Delta}IC 全部为负 \ensuremath{-}0.003 \textasciitilde{} \ensuremath{-}0.013） \\
\bottomrule
\end{longtable}
}
\textbf{降维真正要替代的是后者。} 若目标是"少参数"，那么正确的用法不是"块等权"，而是：
\begin{quote}
\textbf{中立压缩（保分不变）\ensuremath{\rightarrow} 在 G 个块上做加权估计}
\end{quote}
这样下游要估的参数是 \textbf{G 个}（G=10 时 55）而不是 \textbf{24 个}（300），\textbf{降维负责减列、加权交给块级估计} ------ 这才是把 \S{}3.2 的收益兑现出来的做法。
\subsection{\texorpdfstring{七、F 检验（按 ETX5460 项目 2 / HAR 的口径）}{七、F 检验（按 ETX5460 项目 2 / HAR 的口径）}}
按你们习惯的\textbf{参数限制检验}（HAR \ensuremath{\subset} AR(10) 那条的同一形式）：
\[
F=\frac{(RSS_r-RSS_u)/q}{RSS_u/(N-k)},\qquad q=J-G=24-G
\]
{\small
\begin{longtable}{p{2.56cm}p{2.13cm}p{5.97cm}p{2.13cm}p{3.41cm}}
\toprule
G & q & 残差平方和涨幅 & F & p \\
\midrule
15 & 9 & +1.58\% & 7256 & < 1e-300 \\
\textbf{10} & 14 & \textbf{+3.20\%} & 9450 & < 1e-300 \\
6 & 18 & +5.49\% & 12620 & < 1e-300 \\
2 & 22 & +8.38\% & 15740 & < 1e-300 \\
\bottomrule
\end{longtable}
}
\textbf{R\ensuremath{^2} 基线 10.79\%，G=10 时降到 10.45\%（绝对 \ensuremath{-}0.34 个百分点）。}
\begin{quote}
\textbf{\textbf{[注意]} 这个检验必须谨慎解读}：残差自由度 \textbf{413 万}，功效极大，\textbf{任何非零损失都会拒绝}。

作为对照，HAR 那条能得出 p = 0.72（不拒绝），\textbf{唯一原因是它只有 1033 个观测}。

所以 \textbf{F 检验回答的是"约束是否严格成立"（答案是不严格），回答不了"损失是否可忽略"} ------ 后者要靠 \S{}3.3 的等价性检验。
\end{quote}
\subsection{\texorpdfstring{八、下一步}{八、下一步}}
{\small
\begin{longtable}{p{1.27cm}p{7.99cm}p{6.95cm}}
\toprule
步骤 & 做什么 & 目的 \\
\midrule
\textbf{1} & \textbf{中立压缩 + 块级 GLS} & 把"减列"的收益兑现成"参数从 300 降到 55" \\
\textbf{2} & 换聚类判据（按 ICIR 份额分块，允许块大小不均） & 检验"分组方式"能否真正带来增量 \\
\textbf{3} & 训练/验证窗互换重跑 & 排除"结论是某一段行情的产物" \\
\textbf{4} & 剔除反彩票族后重跑压缩曲线 & 确认价值/质量族能独立支撑 \\
\textbf{5} & \textbf{冻结样本外，一次性最终检验} & 到此为止没有碰过样本外数据 \\
\bottomrule
\end{longtable}
}
\textbf{三条贯穿全程的硬规矩}：(1) 所有估计只用训练窗口 (2) 不做"试到通过"的搜索 (3) 每个方法都要说清能做什么、不能做什么。
\subsection{\texorpdfstring{附录 A \textperiodcentered{} 标注汇总}{附录 A · 标注汇总}}
{\small
\begin{longtable}{p{8.10cm}p{8.10cm}}
\toprule
内容 & 标注 \\
\midrule
F 检验 / 参数限制检验 & 🟢 课内（ETX5460 W9 + Project 2 Q3，HAR \ensuremath{\subset} AR(10)） \\
HAC / Newey-West 稳健标准误 & 🟢 课内（ETX5460 W9） \\
层次聚类 / Ward 判据 & 🔴 发散（机器学习） \\
因子构造（波动率 / 反转 / 偏度 / MAX / 非流动性...） & 🔴 发散（业界通用做法） \\
等价性检验（TOST）/ 非劣性检验 & 🔴 发散（临床试验统计的标准工具） \\
安慰剂 / 随机分组对照 / 有效因子数 PR & 🔴 发散（诊断方法） \\
\bottomrule
\end{longtable}
}
\subsection{\texorpdfstring{附录 B \textperiodcentered{} 交付物}{附录 B · 交付物}}
{\small
\begin{longtable}{p{3.24cm}p{12.96cm}}
\toprule
类型 & 文件 \\
\midrule
诊断脚本 & \texttt{\_diag\_dim.R}（降维可行性诊断） \\
扩池脚本 & \texttt{\_factor\_ext2.R}（12 \ensuremath{\rightarrow} 24 因子） \\
字段语义探查 & \texttt{\_probe\_units.R}（价格 / 市值 / 换手的口径校验） \\
聚类脚本 & \texttt{\_stage1\_cluster.R}（两个相似度基准对照） \\
基准对照 & \texttt{\_stage2b\_baseline.R}（6 个变体 + 配对检验） \\
压缩曲线 & \texttt{\_stage3\_compress.R} \\
检验脚本 & \texttt{\_stage4\_equiv.R}（F 检验 + 等价性检验） \\
质疑脚本 & \texttt{\_stage5\_challenge.R}（6 类对抗性检验） \\
数据产物 & \texttt{\_factor\_panel2.parquet}（24 因子面板，1.84 GB） \\
\bottomrule
\end{longtable}
}
\begin{quote}
文档信息：\texttt{Desktop\textbackslash{}就业\textbackslash{}东方证券alpha优化\textbackslash{}因子降维\_质疑与结论.md} ｜生成 2026-10-11 ｜源 \texttt{.md} / \texttt{.tex} 同目录，\textbf{PDF 为交付物}
\end{quote}
\end{document}
